% !PRINTED 649
\weberpage{649}
%<<CONT
so ist $\mathfrak{G}\,|\,\mathfrak{A}$ (\S.\,4 dieses Bandes) die Gruppe des Körpers $\Omega_1$ und
dies ist zugleich mit $\mathfrak{G}$ eine \gesperrt{Abel'sche} Gruppe. Ein Theiler
von $\Omega$, der mit $\Omega$ von gleichem Grade ist, ist mit $\Omega$ identisch
(vgl. Bd. I, \S.\,144, 149, 156).

Giebt es nun zwei nicht primitive Zahlen $x_1, x_2$ in $\Omega$ von
der Art, dass
\[
\Omega = R(x_1, x_2)
\]
gesetzt werden kann, so sind die Körper $\Omega_1 = R(x_1)$, $\Omega_2 = R(x_2)$
echte Theiler von $\Omega$, und $\Omega$ heisst aus den beiden Körpern
$\Omega_1, \Omega_2$ \gesperrt{zusammengesetzt}.

Gestattet der Körper $\Omega$ keine solche Darstellung, so heisst
er \gesperrt{einfach}.

Wenn also als bewiesen vorausgesetzt wird, dass $\Omega_1, \Omega_2$
Kreistheilungskörper sind, d.\,h. dass alle ihre Zahlen rational
durch Einheitswurzeln darstellbar sind, so folgt dasselbe für $\Omega$.

Wenn einer der Körper $\Omega_1, \Omega_2$ zerlegbar ist, so können
wir die Zerlegung wiederholen, und müssen, da die Grade abnehmende
 ganze positive Zahlen sind, endlich zu einfachen Körpern
 gelangen.

\gesperrt{Das Theorem I. braucht also nur für einfache Abel'sche
Körper bewiesen zu werden.}

Es ist aber daran zu erinnern, dass die Zerlegbarkeit eines
Körpers $\Omega$ keineswegs schon aus der Existenz eines Theilers
folgt, und dass bei einem zerlegbaren Körper die Zerlegung in
einfache Körper auf ganz verschiedene Arten geschehen kann, so
dass bei den Körpern nicht die Gesetze der Theilbarkeit gelten,
wie im Gebiete der Zahlen.

Ein Kennzeichen für die Einfachheit eines Körpers können
wir aus seiner Gruppe ableiten.

Wenn die Gruppe $\mathfrak{G}$ zwei echte Theiler $\mathfrak{A}$ und $\mathfrak{B}$ von der
Art hat, dass jedes Element ein- und nur einmal aus einem Element
 von $\mathfrak{A}$ und einem Element von $\mathfrak{B}$ zusammengesetzt werden
kann, so nennen wir die Gruppe $\mathfrak{G}$ zerlegbar in die beiden
Componenten $\mathfrak{A}, \mathfrak{B}$; $\mathfrak{G}$ ist also zerlegbar, wenn in dem nach der
Composition der Theile (\S.\,4) gebildeten Producte
\weq{1}{\mathfrak{G} = \mathfrak{A}\mathfrak{B}}
jedes Element von $\mathfrak{G}$ ein- und nur einmal erscheint. Der Grad
von $\mathfrak{G}$ ist dann gleich dem Producte der Grade von $\mathfrak{A}$ und $\mathfrak{B}$.

Giebt es solche Theiler nicht, so heisst die Gruppe $\mathfrak{G}$
unzerlegbar.
